% !Mode:: "TeX:UTF-8"
\chapter{相关工作与研究现状}

\section{多模态感知}

在多模态感知领域，国际研究前沿已率先利用 CLIP、LLaVA 以及 GPT-4o 等视觉语言大模型（VLM）实现了技术的代际演进，其核心贡献在于通过大规模弱监督学习构建了视觉与文本的统一表征空间。这种进步使得研究重心从单纯的物体检测识别，转向了对复杂驾驶场景的深层语义解析与逻辑理解，尤其在解释自动驾驶中的长尾场景——如极端天气下的模糊障碍物判别或非常规驾驶行为的意图推断方面，取得了显著的理论突破。相比之下，国内研究则紧密结合智慧城市与平安交通的实际业务需求，侧重于利用大模型对海量监控影像进行自动化的时空特征提取，并结合自然语言处理技术生成标准化的事故结构化报告，极大提升了非结构化交通数据在管理链路中的处理效率。

然而，当前国内外研究在迈向全场景感知时共同面临着深层的技术瓶颈：如何实现激光雷达点云的高精度、结构化物理空间信息与视频流中丰富的视觉语义信息在特征层面的深度对齐。这种异构数据融合不仅面临着点云的稀疏性与像素的密集性导致的计算维度不匹配问题，还涉及到跨模态特征在隐空间映射过程中的信息丢失难题。此外，如何在庞大的大模型推理框架下实现毫秒级的实时协同感知，以支撑瞬息万变、对安全性要求极高的交通运行环境，仍是目前全球学术界与工业界亟待攻克的关键技术瓶颈。现有的研究虽然在离线分析上表现卓越，但在构建具备实时反馈能力的“感知-认知”闭环系统方面，依然缺乏兼顾精度与算力开销的优化范式。

\section{交通态势预测}

交通态势预测的技术范式正经历着从传统随机过程或浅层机器学习模型，向以 BigCity 与 LibCity 为代表的时空基座模型的重大转变。此类前沿研究充分利用图卷积网络（GCN）捕捉复杂的路网空间拓扑关联，并结合 Transformer 架构的长效注意力机制来提取多尺度的时间周期特征，从而在城市级大规模交通流预测中展现出了卓越的非线性建模能力。尽管国内科研团队在大规模路网的时空依赖性建模、处理超高维度异构流量数据方面具备显著的工程适配与算法优化优势，但现有的时空模型在应对突发性异常事件（如严重的交通事故、非法占路施工或极端气象灾害）时，依然表现出明显的预测滞后性与鲁棒性不足。

其深层根源在于，目前的定量流量预测模型本质上是基于历史数值序列的概率映射，难以实时、有效地吸收并内化由图像语义捕捉到的定性突发信息，例如视频监控中识别出的车辆侧翻、路面散落物或突发的交通管制指令。这种“数值逻辑”与“视觉语义”在特征层面的严重脱节，导致模型在面临非周期性的极端波动情境时，无法通过因果逻辑修正预测轨迹，其预测精度受限于历史数据的分布，难以实现真正意义上的全场景智能预警。因此，如何构建一种能够融合动态语义感知与稳健时空推理的新型架构，使模型具备对突发扰动的即时理解与演化预判能力，已成为当前交通时空大数据挖掘领域亟待突破的核心难题。

\section{交通智能体}

在智能体（Agent）架构的研究中，大语言模型（LLM）正逐渐演变为整合感知（Perception）、记忆（Memory）、规划（Planning）与执行（Execution）的核心中枢。国际学术界目前侧重于探索 Agent 如何通过调用外部插件与交通仿真引擎（如 SUMO 或 CityFlow），来处理高维度的复杂路径规划与城市级交通流分配问题。与此同时，国内研究领域也涌现出以 TrafficGPT 为代表的垂直领域模型，尝试通过自然语言的交互式指令拆解，驱动底层算法工具链——包括时空基座模型与视觉感知算子——进行自动化的协同调度与逻辑推演。这种“大脑+工具链”的范式初步验证了利用大模型大规模预训练知识来弥补传统交通模型逻辑缺失的可行性。

然而，通用大模型在交通这一专业性极强且容错率极低的垂直场景下，其闭环推理的实际可行性仍面临巨大挑战。尤其是在细粒度的管理意图拆解、严苛的交通法律法规约束以及毫秒级实时响应的可执行性方面，现有智能体架构的专业知识密度与因果推理精度仍显不足。目前的交通智能体在面对复杂的交通事故现场处置建议时，往往容易产生不符合物理常识或违反交通管理规程的“幻觉”输出，且难以在复杂的动态路网环境中实现感知数据与决策动作的实时双向闭环。这深刻表明，如何构建一个深度集成时空感知能力、具备检索增强生成（RAG）支撑的专业知识库、且受到交通领域垂直知识严格约束的专用智能体，已成为当前全球智慧交通与人工智能交叉研究领域的最前沿趋势。
